{
\clearpage
\phantomsection
\addcontentsline{toc}{sectionnoiddot}{《本科生毕业论文（设计）任务书》}
}

\begin{center}
  \setfont{-2}{1.5}\bfseries 本科生毕业论文（设计）任务书
\end{center}

\blankline{5}{1}

\noindent {\setfont{-4}{1.5}\bfseries 一、题目：高能效高精度全动态放大器研究}\par
\noindent {\setfont{-4}{1.5}\bfseries 二、指导教师对毕业论文（设计）的进度安排及任务要求}\par
%----------正文----------

任务要求：

（1） 熟悉动态放大器的研究背景、高能效放大器在混合信号集成电路的应用价值。

（2） 学习电路设计工具，并熟悉Cadence等EDA工具相关设计环境。

（3） 研究动态放大器的关键技术，通过优化电路在增益、带宽、噪声和稳定性之间寻找设计平衡。

（4） 基于标准CMOS工艺设计全动态放大器架构，实现功耗、带宽灵活可调，同时实现低失调和高增益，以满足不同物联网应用场景。

（5） 根据新的电路设计进行仿真实验，进行数据采集与比较分析。

（6） 完成毕业论文（设计）的总结工作和毕业论文的撰写（11000字以上）。



进度安排（以下时间分段供参考，建议按一周或两周间隔进行安排）：

2025-11-10至2025-11-16：明确毕设任务、目标和进度要求。

2025-11-17至2025-11-30：查阅文献资料，调研国内外研究现状，撰写文献综述文档。

完成外文翻译，撰写文献翻译文档。

2025-12-1至2024-12-14：在导师指导下开展实质研究或设计工作。

2025-12-15至2025-12-21：撰写开题报告，应包含学生本人已开展的阶段性工作及结果。

2025-12-22至2025-12-29：讨论修改文献综述、外文翻译和开题报告，准备开题答辩。

2026-01-05至2026-01-11：开题答辩。

2026-03-02至2026-03-22：继续开展毕设研究和/或设计工作。

2026-03-23至2026-03-27：中期检查。

2026-03-28至2026-04-12：继续毕设研究工作，实现毕设研究目标。

2026-04-13至2026-04-20：撰写、上传毕业论文。

2026-04-21至2026-05-24：答辩前审核流程。

2026-05-25至2026-05-29：论文答辩。


%----------正文----------
\vfill
{\setfont{-4}{1.5}\bfseries
\noindent 起讫日期\hspace{1em}2026年 11月 10日至 2026年 05月 29日\vspace{0.5ex}\par
\hfill 指导教师（签名）{\underline{\makebox[7.5em][c]{}}}\hspace{0.5em}
职称{\underline{\makebox[6.5em][c]{教授}}}\vspace{0.5ex}
}

\noindent {\setfont{-4}{1.5}\bfseries 三、系或研究所审核意见}\par
\multido{}{3}{\blankline{-4}{1.5}}
{\setfont{-4}{1.5}\bfseries
\hfill 负责人（签名）{\underline{\makebox[8em][c]{}}}\par
\hfill 2026年XX月XX日
}